\section{Graded generator functors and an explicit Frobenius relation cone}
\label{sec:fgb}

We retain the four actual source classes in their fixed order
\(e_0=T,e_1=T_1,e_2=T_2,e_3=Q\), with precisely the source rescalings
and column words proved in the generator-interval calculation. Put
\[
 \begin{gathered}
 A=\mathbb Z[q,q^{-1}],\quad
 \alpha=[7]-[3],\quad\beta=[8],\quad\gamma=[6],\quad d=[3],\\
 \delta=-d,\quad p=[2],\quad C=q^{10}-q^4-q^{-4}+q^{-10}.
 \end{gathered}
\]
The source quotient has lowering matrix
\[
 N=\begin{pmatrix}0&0&0&0\\\alpha&0&0&0\\
 \beta&0&0&0\\0&\gamma&-d&0\end{pmatrix},\qquad
 D=CE_{30},\qquad N^2=pD.
\]
The construction below gives actual additive functors defined by tensor
products, an integral map
between their two source-labelled length-two branches, and a deformation
retraction after two-periodization to the previously computed Frobenius
object. Its relation map is specified explicitly; it is not inferred from
the equality of Laurent polynomials.

\subsection{The state category and its exact integral decategorification}

Let \(\mathcal G\) be the category of finite free internally
\(\mathbb Z\)-graded abelian groups, each furnished with a homogeneous
basis. Write \(\mathbb Z\langle w\rangle\) for a copy of \(\mathbb Z\)
whose basis vector has internal degree \(w\). Morphisms preserve internal
degree. Let \(\mathcal K=K^b(\mathcal G)\) be its bounded homotopy category,
and let \(\mathcal S=\mathcal K^4\). The tensor differential convention is
\[
 d(x\otimes y)=dx\otimes y+(-1)^i x\otimes dy
 \quad(x\text{ of cohomological degree }i).
\]
Internal degrees add under tensor product. The internal shift
\(\langle w\rangle\) increases every internal degree by \(w\).
The cochain translation convention is \(V[1]^i=V^{i+1}\), with
translated differential \(-d\); in particular a degree-zero group
translated by \([1]\) lies in degree \(-1\).

Define the decategorification group using isomorphism classes in
\(\mathcal K\), direct-sum relations, and
\([\operatorname{Cone}(f)]=[Y]-[X]\) for every cochain map
\(f:X\to Y\). Its internal shift gives an \(A\)-module structure.
There is an isomorphism
\[
 \chi_q:K_0(\mathcal K)\longrightarrow A,
 \qquad [V]\longmapsto
 \sum_{i,w}(-1)^i\operatorname{rank}_{\mathbb Z}(V^i_w)q^w.
 \tag{*}
\]
Here and below sums are finite. To prove the assertion, a direct sum adds
ranks, and \(\operatorname{Cone}(f)^i=Y^i\oplus X^{i+1}\) gives exactly
the required subtraction. A contractible complex has Euler rank zero
in each internal degree: its contracting homotopy splits its successive
boundaries, so the adjacent boundary contributions cancel. Alternatively,
the cone of a homotopy equivalence is contractible, and its finite
alternating chain ranks vanish by the same splitting. Thus the formula
is unchanged under isomorphism in the homotopy category. Filtering a
bounded complex by its cohomological degrees and using the cone relation
expresses its class as the alternating sum of the classes of its chain
groups. Each such group is a direct sum of the specified
\(\mathbb Z\langle w\rangle\). This proves that the map
\(q^w\mapsto[\mathbb Z\langle w\rangle]\) is inverse to
\((\ast)\). The tensor formula and the displayed tensor differential
show that \(\chi_q\) also preserves products. The same argument in the
four factors gives a fixed isomorphism
\[
 \Psi:K_0(\mathcal S)\longrightarrow M_A=\bigoplus_{i=0}^3Ae_i,
 \qquad[(V_0,V_1,V_2,V_3)]\longmapsto
 \sum_{i=0}^3\chi_q(V_i)e_i.
\]
In particular the four labels have not been replaced or permuted.

\subsection{Objects on the four original arrows}

For each indicated finite set, retain the displayed label as part of
the basis. Define degree-zero complexes
\[
 \begin{aligned}
 \mathsf A&=\bigoplus_{a\in\{-6,-4,4,6\}}\mathbb Z a,
       &\deg_q(a)&=a,\\
 \mathsf B&=\bigoplus_{b\in\{-7,-5,-3,-1,1,3,5,7\}}\mathbb Z b,
       &\deg_q(b)&=b,\\
 \mathsf G&=\bigoplus_{g\in\{-5,-3,-1,1,3,5\}}\mathbb Z g,
       &\deg_q(g)&=g,\\
 \mathsf J&=\bigoplus_{j\in\{-2,0,2\}}\mathbb Z j,
       &\deg_q(j)&=j.
 \end{aligned}
\]
The equality \(\chi_q(\mathsf A)=[7]-[3]\) follows by retaining
the seven terms of \([7]\) and subtracting its three terms with
exponents \(2,0,-2\); the four surviving terms are exactly the four
displayed summands. This equality does not rescale any source class.
The other three positive characters are \([8],[6],[3]\), by their
complete finite sums. The negative source sign is implemented by the
cochain translation \(\mathsf J[1]\), which has character \(-[3]\).

Define an exact additive functor \(\mathsf F:\mathcal S\to\mathcal S\)
using tensor products by
\[
 \mathsf F(V)_0=0,\qquad
 \mathsf F(V)_1=\mathsf A\otimes V_0,\qquad
 \mathsf F(V)_2=\mathsf B\otimes V_0,\qquad
 \mathsf F(V)_3=(\mathsf G\otimes V_1)
                       \oplus(\mathsf J[1]\otimes V_2).
\]
Its action on a tuple of cochain maps is the corresponding identity
tensor those maps. Tensoring by these finite free complexes preserves
mapping cones, with the usual tensor-differential signs, so it induces
an endomorphism of the stated decategorification group. Direct use of
the preceding formulas proves the intertwining identity
\[
                 \Psi\,[\mathsf F]=N\,\Psi.
\]
This is an integral identity over the original ring \(A\).

Let \(\mathsf K_V\) shift the four components internally by
\((9,7,7,5)\), respectively, and let \(\mathsf K_W\) shift each
component by \(3\). Their inverses are the opposite shifts. These are
actual invertible functors. For every nonzero edge of \(\mathsf F\),
the target \(V\)-shift is the source shift minus two. Identifying equal
internal shifts on each tensor factor therefore gives a natural
isomorphism
\[
 \mathsf K_V\mathsf F\mathsf K_V^{-1}
                  \simeq\mathsf F\langle-2\rangle.
\]
The \(W\)-shifts are equal on both ends of each edge and commute
with \(\mathsf F\). Set \(\mathsf F_W=0\). Its commutation and
Cartan relations hold as equal zero functors. The two Cartan functors
commute componentwise. Each of the two original central total-weight
generators is represented by the common internal shift \(\langle15\rangle\)
on all four components; these commute with every displayed functor.
Finally \(\mathsf F^3=0\): the only length-two
paths end in component three, from which there is no outgoing arrow.
Thus the negative-Borel relations involving the first lowering
generators and Cartan operators hold already as specified natural
isomorphisms, not only as scalar equalities.

\subsection{The retained geometric scalar object}

Let \(\mathsf C_X\) be the zero-differential integral bigraded complex
with exactly four basis vectors
\[
\begin{array}{c|rrrr}
\text{label}&x_{00}&x_{03}&x_{70}&x_{73}\\\hline
\text{cohomological degree}&2&7&15&20\\
\text{internal degree}&-10&-4&4&10
\end{array}
\]
These labels retain, in order, the K\"unneth factors
\((H^1_c(U_7),H^1_c(U_3))\),
\((H^1_c(U_7),H^6_c(U_3))\),
\((H^{14}_c(U_7),H^1_c(U_3))\), and
\((H^{14}_c(U_7),H^6_c(U_3))\).
The actual groups and their localization and trace identifications
were computed for
\(X=(\mathbb A^7\setminus0)\times(\mathbb A^3\setminus0)\).
Their Tate twists after the common twist \((5)\) are respectively
\((5),(2),(-2),(-5)\). Thus the internal degrees retain twice the
Frobenius exponents and \(\chi_q(\mathsf C_X)=C\), with every
cohomological degree and sign retained.

Define \(\mathsf D(V)\) to have only component three nonzero, equal
to \(\mathsf C_X\otimes V_0\). Then
\(\Psi[\mathsf D]=D\Psi\). Its Cartan relation uses the original
weight difference \(-4\). The support also gives
\(\mathsf F\mathsf D=\mathsf D\mathsf F=\mathsf D^2=0\).
Put \(\mathsf H=\mathbb Z\langle1\rangle\oplus
\mathbb Z\langle-1\rangle\) in cohomological degree zero, so that
\(\chi_q(\mathsf H)=p\).

The following obstruction is exact for these specified objects.
On the state \((\mathbb Z,0,0,0)\), the composite
\(\mathsf F^2\) has component-three complex
\[
  (\mathsf G\otimes\mathsf A)
     \oplus(\mathsf J\otimes\mathsf B)[1],
\]
whose zero differential leaves 24 even and 24 odd free generators.
The complex \(\mathsf H\otimes\mathsf C_X\) has four even and four
odd generators. Tensoring with \(\mathbb Q\), an isomorphism in either
the bounded homotopy category or the two-periodic derived category
would induce an isomorphism of its cohomology. The indicated dimensions
are unequal. Therefore these specified functors have no natural
isomorphism \(\mathsf F^2\simeq\mathsf H\otimes\mathsf D\).
This obstruction proves the failure of that lift; the next map supplies
an exact stronger correspondence between its two branches and the
geometric endpoint.

\subsection{A complete integral map between the two length-two branches}

Write
\[
 \mathsf P=\mathsf G\otimes\mathsf A,\qquad
 \mathsf R=\mathsf J\otimes\mathsf B.
\]
For basis naming use pairs \((a,g)\) in \(\mathsf P\) and \((b,j)\)
in \(\mathsf R\), preserving their source intermediate states
\(T_1\) and \(T_2\), respectively. At each internal weight \(w\),
list the pairs of either type in ascending lexicographic order and
let their lengths be \(p_w,r_w\). Define
\[
 \theta:\mathsf R\longrightarrow\mathsf P,
 \qquad
 \theta(R_{w,k})=
 \begin{cases}P_{w,k},&1\leq k\leq\min(p_w,r_w),\\
 0,&k>p_w.
 \end{cases}
\]
Here \(R_{w,k}\) denotes the \(k\)-th source pair and \(P_{w,k}\)
the \(k\)-th target pair of the given weight. This defines an actual
matrix over \(\mathbb Z\) consisting of twenty distinct entries one
and all remaining entries zero. It preserves internal and cohomological
degree. Its exact ranks and unmatched vectors are
\[
\begin{array}{c|r|l|l}
w&\operatorname{rank}\theta_w&\text{unmatched }(a,g)
                                      &\text{unmatched }(b,j)\\\hline
11&0&(6,5)&\\
9&1&(6,3)&\\
7&2&&\\
5&2&&(7,-2)\\
3&2&&(5,-2)\\
1&3&&\\
-1&3&&\\
-3&2&&(-1,-2)\\
-5&2&&(-3,-2)\\
-7&2&&\\
-9&1&(-4,-5)&\\
-11&0&(-6,-5)&
\end{array}
\]
For verification, the \(\mathsf P\)-basis is the Cartesian product
of the four explicitly given values of \(a\) and the six values of
\(g\); the \(\mathsf R\)-basis is the product of the eight values
of \(b\) and the three values of \(j\). Adding the two entries gives
exactly the twelve rows in the table. Listing and matching within one
row cannot affect any other row. The rank column sums to twenty,
and each side has 24 basis vectors, so each side has exactly the
four listed residual generators. Because all matched entries are
one on distinct rows and columns, the image is a direct summand and
both the kernel and cokernel are free, with exactly those residual
bases. There is no unrecorded integral torsion.

Let \(\mathsf L=\operatorname{Cone}(\theta)\). It has
\(\mathsf R\) in cohomological degree \(-1\), \(\mathsf P\) in
degree zero, and differential exactly \(\theta\). In this convention
the sign of its character is
\[
 \chi_q(\mathsf L)=\gamma\alpha-d\beta
                 =pC.
\]
More strongly, the complex splits integrally as twenty two-term
complexes \(\mathbb Z\xrightarrow{1}\mathbb Z\) and eight isolated
residual generators. The split is the displayed basis permutation,
not a rational rank argument. Map each matched target basis vector
back to its matched source basis vector and all other basis vectors
to zero; call this cochain homotopy \(h\). If \(\rho\) is projection
to the eight residual vectors and \(\iota\) is their inclusion, then
on a matched pair \(dh+hd\) is the identity on both vectors and on
a residual vector it is zero. Consequently
\[
                    1-\iota\rho=dh+hd,
                    \qquad \rho\iota=1.
\]
This proves an explicit integral deformation retraction and all its
signs.

Regard \(\mathsf P,\mathsf R,\mathsf L\) also as endpoint functors
on \(\mathcal S\), tensoring their arguments from component zero
into component three. The map \(\theta\otimes1\) is natural in
every state and morphism. Its mapping-cone triangle is therefore
a triangle of explicitly specified functors
\[
             \mathsf R\xrightarrow{\theta}\mathsf P
                    \longrightarrow\mathsf L\longrightarrow\mathsf R[1].
\]
The original composite \(\mathsf F^2\) is the split sum
\(\mathsf P\oplus\mathsf R[1]\), not this cone with its nonzero
cross-branch differential. Their equality in the decategorification
group follows from this actual triangle. Their nonisomorphism follows
from the already computed cohomology dimensions. In particular,
the cross-branch relation has not silently been installed as the
tensor differential of \(\mathsf F^2\).

\subsection{The endpoint map with all cohomological degrees recorded}

Let \(\operatorname{Per}\) send a bounded complex to its two-periodic
complex by summing its even and odd cohomological degrees separately,
retaining the differential, internal degree, and individual summand
labels. This operation explicitly forgets the distinction between
cohomological degrees of the same parity. The eight residual
generators of \(\mathsf L\) identify with the eight indicated tensor
generators of \(\mathsf H\otimes\mathsf C_X\) as follows. The symbol
\(+\) means the \(\mathsf H\)-basis of internal degree one and
\(-\) means its basis of internal degree minus one:
\[
\begin{array}{c|c|r|r}
\text{residual path}&\mathsf H\otimes\mathsf C_X
 &\text{internal degree}&\text{original endpoint cohomological degree}\\\hline
(6,5)&+\otimes x_{73}&11&20\\
(6,3)&-\otimes x_{73}&9&20\\
(-4,-5)&+\otimes x_{00}&-9&2\\
(-6,-5)&-\otimes x_{00}&-11&2\\
(7,-2)&+\otimes x_{70}&5&15\\
(5,-2)&-\otimes x_{70}&3&15\\
(-1,-2)&+\otimes x_{03}&-3&7\\
(-3,-2)&-\otimes x_{03}&-5&7
\end{array}
\]
The first four residual vectors are in cohomological degree zero
in \(\mathsf L\), while the last four are in degree \(-1\).
Each row preserves parity and internal degree. Extend this table
linearly to an isomorphism of the residual two-periodic complex
with \(\operatorname{Per}(\mathsf H\otimes\mathsf C_X)\).
Composition with \(\rho\), and in the other direction with
\(\iota\), gives inverse homotopy equivalences
\[
 \operatorname{Per}(\mathsf L)
       \simeq\operatorname{Per}(\mathsf H\otimes\mathsf C_X),
\]
with precisely the contraction \(h\) proved above. This supplies
the actual morphism relating all surviving branch summands to the
scalar geometry. The cohomological columns also prove that this
map is not a degree-preserving map to the unperiodized scalar
cohomology: for example its weight-eleven class has degree zero
on the cone side and degree twenty on the geometric side.
Moreover the differences for the weights eleven and minus eleven
are twenty and two, respectively, so no single common cohomological
shift repairs this failure.

\subsection{Frobenius realization and the integral-Tate parity obstruction}

Fix a prime power \(Q\), a prime \(\ell\) not dividing \(Q\), and
choose a root \(z\) of \(z^2=Q\) in
\(E=\mathbb Q_\ell(z)\). This is a specified extra coefficient
choice for the odd internal degrees; it was unnecessary for the
even-degree scalar object. The element \(z\) is a unit in the
ring of integers of \(E\). For each power of its maximal ideal,
its reduction has finite multiplicative order, hence
\(n\mapsto z^n\) extends from \(\mathbb Z\) to a continuous
homomorphism from \(\widehat{\mathbb Z}\) to that finite unit
group. The extensions commute with reduction because they agree
on the dense integer subgroup. Their inverse limit is a continuous
character \(\widehat{\mathbb Z}\to E^\times\).
Identifying the Galois group of \(\overline{\mathbb F}_Q/\mathbb F_Q\)
with \(\widehat{\mathbb Z}\), with geometric Frobenius sent to one,
therefore defines an actual one-dimensional continuous representation
\(\mathcal L\) on which geometric Frobenius acts as \(z\).
Its square is identified, on their specified basis vectors, with
\(E(-1)\); both characters send geometric Frobenius to \(Q\).

Define the tensor realization on every graded basis summand by
\[
 \mathbb Z\langle w\rangle\longmapsto\mathcal L^{\otimes w},
\]
using dual powers for negative \(w\), extending scalars in the
integer matrices to \(E\), and retaining every cohomological degree.
Maps preserving internal degree become Frobenius-equivariant maps,
because every nonzero matrix entry joins the same character on
source and target. In particular the twenty entries of \(\theta\),
the projection, inclusion and contraction, and every entry of the
residual table are Frobenius equivariant. The actual geometric
object \(H^*_c(X,E)(5)\) is the realization of \(\mathsf C_X\)
by its previously proved localization and K\"unneth identifications.
Here \(H^*_c\) explicitly denotes the direct sum of the actual
cohomology groups placed in their original cohomological degrees with
zero differential. It does not denote a chosen cochain model for
\(R\Gamma_c(X,E)\). No canonical splitting of that derived object
into its cohomology groups is asserted or used.
We obtain a specified equivalence of two-periodic complexes of
continuous Frobenius representations
\[
 \operatorname{Per}(\mathsf L_E)
 \simeq\operatorname{Per}\bigl((\mathcal L\oplus\mathcal L^{-1})
                                  \otimes H^*_c(X,E)(5)\bigr).
\]
Every entry of the map and its inverse has been given above. Its
alternating trace of the \(m\)-th Frobenius power is
\[
 (z^m+z^{-m})(Q^{5m}-Q^{2m}-Q^{-2m}+Q^{-5m}).
\]
If the chosen complex embedding of \(z\) is the positive square
root, this is precisely \([2]_{Q^{m/2}}C(Q^{m/2})\). The two choices
of \(z\) are retained choices of odd-degree realization; the four
even-degree scalar summands are independent of that choice.

The need for this extension has a precise parity obstruction. Suppose
one tries to realize all four arrow coefficients and the divided-square
coefficient by only integral Tate degrees after replacing the four
state generators by \(q^{s_i}e_i\), for integers \(s_i\). A coefficient
from state \(i\) to state \(j\) becomes multiplied by
\(q^{s_i-s_j}\). Integral Tate characters have only even powers of
\(q\). The nonzero exponent parities of \(\alpha,\beta,\gamma,-d\)
are respectively \(0,1,1,0\). Thus all four arrows being even requires
\[
 s_0-s_1=0,\quad 1+s_0-s_2=0,\quad
 1+s_1-s_3=0,\quad s_2-s_3=0\pmod2.
\]
The first and third imply \(s_3-s_0=1\pmod2\). But every term of
\(C\) has even exponent, so the divided-square arrow from zero to
three being even would require \(s_3-s_0=0\pmod2\), a contradiction.
This proves that diagonal state regradings cannot provide that
all-integral-Tate lift. The character \(\mathcal L\) just constructed
is an explicit alternative, not an unmentioned half-integral Tate
twist in the original scalar calculation.

The construction therefore gives an integral four-state generator
functor, all its first-generator negative-Borel relations, a precise
decategorification morphism to the actual source subquotient, a complete
rank-twenty relation map, and its two-periodic Frobenius endpoint
equivalence. It also records the exact obstruction to the naive divided
power functor isomorphism and to restoring the full cohomological
grading by a common shift. It does not identify the added relation map
with a canonical-basis convolution map in the full source module.
Constructing that source geometric compatibility remains a further
mathematical problem; no nonstandard Riemann-hypothesis or complexity
conclusion is supplied by the finite interval alone.
